% GENERATED FILE. Edit canonical lessons, then run npm run generate:books.
% canonical-source-hash: fnv1a64:ae13105af38518fe
% canonical-lessons: AR-C379-hizb, AR-C379-hamla, AR-C379-tarashshaha, AR-C379-nashit, AR-C379-dimuqratiyya

\chapter{Parties and campaigns}
\label{ch:ar-parties-and-campaigns}
\hlchaptermodality{379}

\begin{chapteropening}
\textbf{By the end of this chapter:} \emph{I can talk about a party, a campaign, an activist and democracy, and say that someone ran for office.}

Complete the last lesson of chapter 379: I can talk about a party, a campaign, an activist and democracy, and say that someone ran for office.
\end{chapteropening}

\section[\texorpdfstring{\ar{حزب} (ḥizb)}{حزب (ḥizb)}]{\ar{حزب} (ḥizb) — a political party}

\label{lesson:AR-C379-hizb}



\begin{quote}
Before the new one: say the Arabic for a voter, then the Arabic for a candidate.
\end{quote}

\subsection*{You'll want to know: \ar{حزب}}
\textbf{\ar{حزب}} — \emph{ḥizb} — \textquotedblleft{}a political party\textquotedblright{}.

\begin{quote}
\ar{لكل} \ar{حزب} \ar{برنامج} \ar{مختلف}.
\end{quote}

\begin{quote}
\emph{li-kulli ḥizb barnāmaj mukhtalif.}
\end{quote}

\begin{quote}
Each party has a different programme.
\end{quote}

Plural: \ar{أحزاب}, \emph{aḥzāb}.

\begin{cousinweb}
\textbf{\ar{حزب}} (\emph{ḥizb}) is built on the native Arabic root \emph{ḥ-z-b} (Lane), where it names a band of people who hold together. Modern Arabic gave it the political party.
\end{cousinweb}

\begin{grammarlens}[title={What you've built}]
One more word for giving reasons and opinions.
\end{grammarlens}

\subsection*{Your turn}
\begin{itemize}
\raggedright
  \item \emph{Say it:} \emph{ḥizb}
  \item \emph{Say it:} \emph{ḥizb}, once more
  \item \emph{Say it:} \emph{li-kulli ḥizb barnāmaj mukhtalif.} — Each party has a different programme.
  \item \emph{Recall:} say the Arabic for a voter, then the Arabic for a candidate, then say \emph{ḥizb} again
\end{itemize}

\begin{tcolorbox}[breakable,colback=teal!4,colframe=teal!35!black,title={Before you move on}]
What does \ar{حزب} mean? (\textquotedblleft{}A political party\textquotedblright{}.) Say it once more.
\end{tcolorbox}
\section[\texorpdfstring{\ar{حملة} (ḥamla)}{حملة (ḥamla)}]{\ar{حملة} (ḥamla) — a campaign}

\label{lesson:AR-C379-hamla}



\begin{quote}
Before the new one: say the Arabic for a candidate, then the Arabic for a political party.
\end{quote}

\subsection*{You'll want to know: \ar{حملة}}
\textbf{\ar{حملة}} — \emph{ḥamla} — \textquotedblleft{}a campaign\textquotedblright{}.

\begin{quote}
\ar{بدأت} \ar{الحملة} \ar{قبل} \ar{شهرين}.
\end{quote}

\begin{quote}
\emph{badaʾat al-ḥamla qabla shahrayn.}
\end{quote}

\begin{quote}
The campaign began two months ago.
\end{quote}

\begin{cousinweb}
\textbf{\ar{حملة}} (\emph{ḥamla}) is built on the native Arabic root \emph{ḥ-m-l}, \textquotedblleft{}to carry, to bear\textquotedblright{} (Lane). A \emph{ḥamla} is one charge, one carrying forward: an attack, and so a campaign.
\end{cousinweb}

\begin{grammarlens}[title={What you've built}]
One more word for giving reasons and opinions.
\end{grammarlens}

\subsection*{Your turn}
\begin{itemize}
\raggedright
  \item \emph{Say it:} \emph{ḥamla}
  \item \emph{Say it:} \emph{ḥamla}, once more
  \item \emph{Say it:} \emph{badaʾat al-ḥamla qabla shahrayn.} — The campaign began two months ago.
  \item \emph{Recall:} say the Arabic for a candidate, then the Arabic for a political party, then say \emph{ḥamla} again
\end{itemize}

\begin{tcolorbox}[breakable,colback=teal!4,colframe=teal!35!black,title={Before you move on}]
What does \ar{حملة} mean? (\textquotedblleft{}A campaign\textquotedblright{}.) Say it once more.
\end{tcolorbox}
\section[\texorpdfstring{\ar{ترشح} (tarashshaḥa)}{ترشح (tarashshaḥa)}]{\ar{ترشح} (tarashshaḥa) — to run for office}

\label{lesson:AR-C379-tarashshaha}



\begin{quote}
Before the new one: say the Arabic for a political party, then the Arabic for a campaign.
\end{quote}

\subsection*{You'll want to know: \ar{ترشح}}
\textbf{\ar{ترشح}} — \emph{tarashshaḥa} — \textquotedblleft{}to run for office\textquotedblright{}.

\begin{quote}
\ar{ترشحت} \ar{جارتنا} \ar{في} \ar{الانتخابات}.
\end{quote}

\begin{quote}
\emph{tarashshaḥat jāratunā fī al-intikhābāt.}
\end{quote}

\begin{quote}
Our neighbour ran in the elections.
\end{quote}

Present: \ar{يترشح}, \emph{yatarashshaḥ}.

It takes \ar{لـ} for the post one runs for.

\begin{cousinweb}
\textbf{\ar{ترشح}} (\emph{tarashshaḥa}) is built on the native Arabic root \emph{r-sh-ḥ}, \textquotedblleft{}to sweat, to ooze\textquotedblright{} (Lane). It is Form V, to make oneself ready for a post: to stand as a candidate.
\end{cousinweb}

\begin{grammarlens}[title={What you've built}]
One more word for giving reasons and opinions.
\end{grammarlens}

\subsection*{Your turn}
\begin{itemize}
\raggedright
  \item \emph{Say it:} \emph{tarashshaḥa}
  \item \emph{Say it:} \emph{tarashshaḥa}, once more
  \item \emph{Say it:} \emph{tarashshaḥat jāratunā fī al-intikhābāt.} — Our neighbour ran in the elections.
  \item \emph{Recall:} say the Arabic for a political party, then the Arabic for a campaign, then say \emph{tarashshaḥa} again
\end{itemize}

\begin{tcolorbox}[breakable,colback=teal!4,colframe=teal!35!black,title={Before you move on}]
What does \ar{ترشح} mean? (\textquotedblleft{}To run for office\textquotedblright{}.) Say it once more.
\end{tcolorbox}
\section[\texorpdfstring{\ar{ناشط} (nāshiṭ)}{ناشط (nāshiṭ)}]{\ar{ناشط} (nāshiṭ) — an activist}

\label{lesson:AR-C379-nashit}



\begin{quote}
Before the new one: say the Arabic for a campaign, then the Arabic for to run for office.
\end{quote}

\subsection*{You'll want to know: \ar{ناشط}}
\textbf{\ar{ناشط}} — \emph{nāshiṭ} — \textquotedblleft{}an activist\textquotedblright{}.

\begin{quote}
\ar{تعمل} \ar{الناشطة} \ar{مع} \ar{الأطفال} \ar{الفقراء}.
\end{quote}

\begin{quote}
\emph{taʿmalu an-nāshiṭa maʿa al-aṭfāl al-fuqarāʾ.}
\end{quote}

\begin{quote}
The activist works with poor children.
\end{quote}

\begin{cousinweb}
\textbf{\ar{ناشط}} (\emph{nāshiṭ}) is built on the native Arabic root \emph{n-sh-ṭ}, \textquotedblleft{}to be brisk, lively\textquotedblright{} (Lane). \emph{Nāshiṭ} is the active participle: one who is active, an activist. \ar{نشيط}, active, is from the same root.
\end{cousinweb}

\begin{grammarlens}[title={What you've built}]
One more word for giving reasons and opinions.
\end{grammarlens}

\subsection*{Your turn}
\begin{itemize}
\raggedright
  \item \emph{Say it:} \emph{nāshiṭ}
  \item \emph{Say it:} \emph{nāshiṭ}, once more
  \item \emph{Say it:} \emph{taʿmalu an-nāshiṭa maʿa al-aṭfāl al-fuqarāʾ.} — The activist works with poor children.
  \item \emph{Recall:} say the Arabic for a campaign, then the Arabic for to run for office, then say \emph{nāshiṭ} again
\end{itemize}

\begin{tcolorbox}[breakable,colback=teal!4,colframe=teal!35!black,title={Before you move on}]
What does \ar{ناشط} mean? (\textquotedblleft{}An activist\textquotedblright{}.) Say it once more.
\end{tcolorbox}
\section[\texorpdfstring{\ar{ديمقراطية} (dīmuqrāṭiyya)}{ديمقراطية (dīmuqrāṭiyya)}]{\ar{ديمقراطية} (dīmuqrāṭiyya) — democracy}

\label{lesson:AR-C379-dimuqratiyya}



\begin{quote}
Before the new one: say the Arabic for to run for office, then the Arabic for an activist.
\end{quote}

\subsection*{You'll want to know: \ar{ديمقراطية}}
\textbf{\ar{ديمقراطية}} — \emph{dīmuqrāṭiyya} — \textquotedblleft{}democracy\textquotedblright{}.

\begin{quote}
\ar{يتكلم} \ar{الكتاب} \ar{عن} \ar{تاريخ} \ar{الديمقراطية}.
\end{quote}

\begin{quote}
\emph{yatakallamu al-kitāb ʿan tārīkh ad-dīmuqrāṭiyya.}
\end{quote}

\begin{quote}
The book talks about the history of democracy.
\end{quote}

\begin{cousinweb}
\textbf{\ar{ديمقراطية}} (\emph{dīmuqrāṭiyya}) is a modern loan, usually traced to Greek \emph{dēmokratia}, \textquotedblleft{}rule of the people\textquotedblright{}, through European languages. Arabic added its own abstract ending \emph{-iyya}.
\end{cousinweb}

\begin{grammarlens}[title={What you've built}]
One more word for giving reasons and opinions.
\end{grammarlens}

\subsection*{Your turn}
\begin{itemize}
\raggedright
  \item \emph{Say it:} \emph{dīmuqrāṭiyya}
  \item \emph{Say it:} \emph{dīmuqrāṭiyya}, once more
  \item \emph{Say it:} \emph{yatakallamu al-kitāb ʿan tārīkh ad-dīmuqrāṭiyya.} — The book talks about the history of democracy.
  \item \emph{Recall:} say the Arabic for to run for office, then the Arabic for an activist, then say \emph{dīmuqrāṭiyya} again
\end{itemize}

\begin{tcolorbox}[breakable,colback=teal!4,colframe=teal!35!black,title={Before you move on}]
What does \ar{ديمقراطية} mean? (\textquotedblleft{}Democracy\textquotedblright{}.) Say it once more.
\end{tcolorbox}
